\begin{scanpage}{0021}{4}
In the second case, the ``magnetic'' vector potential $A_z^m$ satisfies the equation
\[\Delta A_z^m+k^2A_z^m=-\frac{4\pi}{c}j_z^m\tag{1.05}\]
and the boundary condition
\[\frac{\partial A_z^m}{\partial\phi}=0\quad\text{with}\quad\phi=0\text{ and }\phi=\alpha.\tag{1.06}\]
It is natural to seek the solution of the nonhomogeneous Equations (1.03) and (1.05) in the form
\[A_z^e=\begin{cases}\displaystyle\sum_{s=1}^{\infty}a_sJ_{\nu_s}(kr)H_{\nu_s}^{(1)}(kr_0)\sin\nu_s\phi_0\sin\nu_s\phi&\text{with }r<r_0,\\\displaystyle\sum_{s=1}^{\infty}a_sJ_{\nu_s}(kr_0)H_{\nu_s}^{(1)}(kr)\sin\nu_s\phi_0\sin\nu_s\phi&\text{with }r>r_0;\end{cases}\tag{1.07}\]
\[A_z^m=\begin{cases}\displaystyle\sum_{s=0}^{\infty}b_sJ_{\nu_s}(kr)H_{\nu_s}^{(1)}(kr_0)\cos\nu_s\phi_0\cos\nu_s\phi&\text{with }r<r_0,\\\displaystyle\sum_{s=0}^{\infty}b_sJ_{\nu_s}(kr_0)H_{\nu_s}^{(1)}(kr)\cos\nu_s\phi_0\cos\nu_s\phi&\text{with }r>r_0,\end{cases}\tag{1.08}\]
\[\nu_s=s\frac{\pi}{\alpha}.\]
The products
\[\begin{array}{ccc}J_{\nu_s}(kr)\sin\nu_s\phi&&J_{\nu_s}(kr)\cos\nu_s\phi\\&\text{and}&\\H_{\nu_s}^{(1)}(kr)\sin\nu_s\phi&&H_{\nu_s}^{(1)}(kr)\cos\nu_s\phi\end{array}\tag{1.09}\]
are the partial solutions of Equations (1.03) and (1.05) without the right-hand member which satisfy the boundary conditions (1.04) and (1.06). The remaining factors entering into Equations (1.07) and (1.08) ensure the observance of the reciprocity principle and the continuity of the field on the arc $r=r_0$. The Bessel function $J_{\nu_s}(kr)$ enters these equations when $r<r_0$, because it remains finite when $r\to0$, and the Hankel function $H_{\nu_s}^{(1)}(kr)$ is taken when $r>r_0$ in order that the solution satisfies the radiation condition.
\end{scanpage}
